\section{Cumulative span}

The exact-word identities used in this section are proved in
Section~\ref{sec:app_wielandt_exact_word_support}. The stabilization and
spectral linear algebra behind the cumulative estimates are collected in
Section~\ref{sec:app_wielandt_cumulative_support}.

\begin{definition}[Word span]\label{def:word_span}
    \lean{MPSTensor.wordSpan}
    \leanok
    \uses{def:mps_tensor, def:eval_word}
    The \emph{word span} at length~$n$ is the subspace
    \begin{align}
        S_n(A)
         & =\spn_{\C}\{A^w:|w|=n\}\subseteq\MN{D}.
        \notag
    \end{align}
    This is~\cite[Section~II]{Sanz2010Quantum}.
\end{definition}

\begin{definition}[Cumulative span]\label{def:cumulative_span}
    \lean{MPSTensor.cumulativeSpan}
    \leanok
    \uses{def:word_span}
    The \emph{cumulative span} at level~$n$ is
    $T_n(A)=S_0(A)+S_1(A)+\cdots+S_n(A)$.
\end{definition}

\begin{definition}[Fixed-length vector span]\label{def:vector_spread_span}
    \lean{MPSTensor.vectorSpreadSpan}
    \leanok
    \uses{def:mps_tensor, def:eval_word}
    The \emph{fixed-length vector span} at length~$n$ is
    \begin{align}
        H_n(A,\varphi)
         & =\spn_{\C}\{A^w\varphi:|w|=n\}\subseteq\C^D.
        \notag
    \end{align}
    This is $S_n(A)\ket{\varphi}$ in the notation
    of~\cite{Sanz2010Quantum}.
\end{definition}

\subsection{Paper primitivity and indices}

For the quantitative bounds, write $d'=\dim S_1(A)$ for the number of
linearly independent Kraus operators. The source definitions and
Proposition~3 of~\cite{Sanz2010Quantum} distinguish uniform vector spreading,
eventual exact-word spanning, and spectral strong irreducibility.

\begin{definition}[Kraus rank]\label{def:wielandt_kraus_rank}
    \lean{MPSTensor.krausRank}
    \leanok
    \uses{def:word_span}
    The \emph{Kraus rank} is
    $\krausrk(A)=\dim S_1(A)$.
\end{definition}

\begin{definition}[Eventually full Kraus rank]
    \label{def:eventually_full_kraus_rank}
    \lean{MPSTensor.HasEventuallyFullKrausRank}
    \leanok
    \uses{def:word_span}
    The Kraus family has \emph{eventually full Kraus rank} if there is a
    positive integer $i$ such that $S_i(A)=\MN{D}$.
\end{definition}

\begin{definition}[Paper primitivity]\label{def:paper_primitive}
    \lean{MPSTensor.IsPrimitivePaper}
    \leanok
    \uses{def:vector_spread_span}
    The transfer map is \emph{primitive in the sense of
        \cite{Sanz2010Quantum}} if there is a positive integer $q$ such that
    \begin{align}
        H_q(A,\varphi)
         & =\C^D
        \label{eq:wld_paper_primitive_spread}
    \end{align}
    for every nonzero $\varphi\in\C^D$. The same $q$ works for all
    nonzero vectors.
\end{definition}

\begin{definition}[Full-Kraus-rank index]\label{def:wielandt_i_index}
    \lean{MPSTensor.iIndex}
    \leanok
    \uses{def:eventually_full_kraus_rank}
    The \emph{full-Kraus-rank index} is
    \begin{align}
        \iota(A)
         & =\min\{n\ge1:S_n(A)=\MN{D}\}.
        \notag
    \end{align}
    If the set is empty, its value is defined to be zero.
\end{definition}

\begin{definition}[Primitivity index]\label{def:wielandt_q_index}
    \lean{MPSTensor.qIndex}
    \leanok
    \uses{def:paper_primitive}
    The \emph{primitivity index} is
    \begin{align}
        q(\E_A)
         & =\min\{q\ge1:H_q(A,\varphi)=\C^D
        \text{ for every }\varphi\neq0\}.
        \notag
    \end{align}
    If the set is empty, its value is defined to be zero.
\end{definition}

% Scope restriction (explicit irreducibility): the condition below strengthens
% Sanz et al., Proposition 3(c). See
% docs/paper-gaps/quantum_wielandt_deviation_v1.tex.
\begin{definition}[Strengthened strong irreducibility]
    \label{def:strongly_irreducible_paper}
    \lean{MPSTensor.IsStronglyIrreduciblePaper}
    \leanok
    \uses{def:transfer_map, def:irreducible_cp, def:primitive_channel}
    We call the tensor \emph{strongly irreducible} here if there is a matrix
    $\rho>0$ such that
    \begin{align}
        \E_A(\rho)                 & =\rho, \notag \\
        \spec_{\mathrm{per}}(\E_A) & =\{1\},
        \notag
    \end{align}
    and $\E_A$ is irreducible. The condition in
    \cite[Proposition~3(c)]{Sanz2010Quantum} requires that $1$ be the only
    peripheral eigenvalue and that its eigenspace be one-dimensional, generated
    by a positive-definite fixed point. The definition above strengthens this
    condition by also requiring irreducibility explicitly.
\end{definition}

The next seven results use the following notation. For nonnegative integers
$d$ and $D$, let $K=\{K^i\}_{i=0}^{d-1}$ be a finite family of $D\times D$
matrices, and let $\E_K$ be its Kraus map.

\begin{theorem}[Iterated Kraus map on a rank-one matrix]
    \label{thm:kraus_iterate_rank_one_sum}
    \lean{Kraus.mapLM_pow_rankOne_eq_sum}
    \leanok
    \uses{def:kraus_map}
    For every choice of $d$, $D$, and $K$ as above, every $q\ge0$, and every
    $\varphi\in\C^D$,
    \begin{align}
        \E_K^q(\ketbra{\varphi}{\varphi})
        =\sum_{|w|=q}\ketbra{K^w\varphi}{K^w\varphi}.
        \notag
    \end{align}
\end{theorem}

\begin{proof}\leanok
    \uses{thm:kraus_map_iterate_word_expansion}
    Expand the Kraus representation of the $q$-fold iterate. Each summand is
    $K^w\ketbra{\varphi}{\varphi}(K^w)^\dagger
      =\ketbra{K^w\varphi}{K^w\varphi}$.
\end{proof}

\begin{theorem}[Positive rank-one image from vector spreading]
    \label{thm:kraus_rank_one_posdef_from_spreading}
    \lean{Kraus.mapLM_pow_rankOne_posDef_of_vectorSpreadSpan_eq_top}
    \leanok
    \uses{def:kraus_map, def:vector_spread_span}
    For every choice of $d$, $D$, and $K$ as above and every $q\ge0$, suppose
    $H_q(K,\varphi)=\C^D$ for every nonzero
    $\varphi\in\C^D$. Then
    $\E_K^q(\ketbra{\varphi}{\varphi})>0$ for every
    $\varphi\neq0$.
\end{theorem}

\begin{proof}\leanok
    \uses{thm:kraus_iterate_rank_one_sum, thm:finite_frame_posdef}
    The preceding theorem expresses the image as the frame operator of the
    spanning family $(K^w\varphi)_{|w|=q}$. A finite frame operator is positive
    definite exactly when its vectors span the ambient space.
\end{proof}

\begin{theorem}[Positive-semidefinite preservation by Kraus iterates]
    \label{thm:kraus_iterate_psd}
    \lean{Kraus.mapLM_pow_posSemidef}
    \leanok
    \uses{def:kraus_map}
    For every choice of $d$, $D$, and $K$ as above and every $n\ge0$, the
    iterate $\E_K^n$ maps positive-semidefinite
    matrices to positive-semidefinite matrices.
\end{theorem}

\begin{proof}\leanok
    \uses{thm:kraus_iterate_rank_one_sum}
    Decompose a positive-semidefinite matrix into a finite sum of rank-one
    positive matrices and apply Theorem~\ref{thm:kraus_iterate_rank_one_sum}
    to every summand.
\end{proof}

\begin{theorem}[Positivity improvement from vector spreading]
    \label{thm:kraus_positivity_improving_from_spreading}
    \lean{Kraus.mapLM_pow_positivityImproving_of_vectorSpreadSpan_eq_top}
    \leanok
    \uses{def:kraus_map, def:vector_spread_span}
    For every choice of $d$, $D$, and $K$ as above and every $q\ge0$, suppose
    $H_q(K,\varphi)=\C^D$ for every nonzero $\varphi$.
    Then $\E_K^q(\rho)>0$ for every nonzero positive-semidefinite matrix
    $\rho$.
\end{theorem}

\begin{proof}\leanok
    \uses{thm:kraus_rank_one_posdef_from_spreading, thm:kraus_iterate_psd}
    Write $\rho$ as a sum of rank-one positive matrices. At least one vector
    in this decomposition is nonzero, so its image is positive definite by
    Theorem~\ref{thm:kraus_rank_one_posdef_from_spreading}; all remaining
    images are positive semidefinite by Theorem~\ref{thm:kraus_iterate_psd}.
\end{proof}

\begin{theorem}[Definiteness of fixed points from vector spreading]
    \label{thm:kraus_fixed_point_posdef_from_spreading}
    \lean{Kraus.posDef_fixedPoint_of_vectorSpreadSpan_eq_top}
    \leanok
    \uses{def:kraus_map, def:vector_spread_span}
    For every choice of $d$, $D$, and $K$ as above and every $q\ge0$, suppose
    $H_q(K,\varphi)=\C^D$ for every nonzero $\varphi$. Then every nonzero
    positive-semidefinite fixed point of $\E_K$ is positive definite.
\end{theorem}

\begin{proof}\leanok
    \uses{thm:kraus_positivity_improving_from_spreading}
    A fixed point of $\E_K$ is fixed by $\E_K^q$, whose action on every
    nonzero positive-semidefinite matrix is positive definite.
\end{proof}

\begin{theorem}[Definiteness of fixed points of positive powers]
    \label{thm:kraus_power_fixed_point_posdef_from_spreading}
    \lean{Kraus.posDef_pow_fixedPoint_of_vectorSpreadSpan_eq_top}
    \leanok
    \uses{def:kraus_map, def:vector_spread_span}
    For every choice of $d$, $D$, and $K$ as above and every $q\ge0$, suppose
    $H_q(K,\varphi)=\C^D$ for every nonzero $\varphi$. For every $p>0$, each
    nonzero
    positive-semidefinite fixed point of $\E_K^p$ is positive definite.
\end{theorem}

\begin{proof}\leanok
    \uses{thm:kraus_positivity_improving_from_spreading, thm:kraus_iterate_psd}
    If $\E_K^p(\rho)=\rho$, then $\E_K^{pq}(\rho)=\rho$. Factor this iterate
    as $\E_K^q\E_K^{(p-1)q}$. The inner image is nonzero and positive
    semidefinite, so the outer application is positive definite.
\end{proof}

\begin{theorem}[Irreducibility from vector spreading]
    \label{thm:kraus_irreducible_from_spreading}
    \lean{Kraus.isIrreducibleMap_mapLM_of_vectorSpreadSpan_eq_top}
    \leanok
    \uses{def:kraus_map, def:vector_spread_span, def:irreducible_cp}
    For every choice of $d$, $D$, and $K$ as above and every $q\ge0$, if
    $H_q(K,\varphi)=\C^D$ for every nonzero $\varphi$, then the Kraus map
    $\E_K$ is irreducible.
\end{theorem}

\begin{proof}\leanok
    \uses{thm:kraus_invariant_compression_lower_zero}
    Let $P$ be an invariant orthogonal projection. Every word $K^w$ preserves
    $\operatorname{ran}P$. If $P\neq0$, choose a nonzero vector in this range.
    Its length-$q$ word orbit spans $\C^D$ by hypothesis and remains in
    $\operatorname{ran}P$, hence $P=\Id$.
\end{proof}

\begin{theorem}[Conjugate eigenmatrix of a positive map]
    \label{thm:positive-map-conj-eigenvector}
    \lean{Kraus.conjTranspose_eigenvector}
    \leanok
    \uses{def:positive_map}
    Let $E:\MN{D}\to\MN{D}$ be positive. If $E(X)=\mu X$, then
    $E(X^\dagger)=\overline\mu X^\dagger$.
\end{theorem}

\begin{proof}\leanok
    \uses{thm:positive_map_conjTranspose}
    Positivity implies that $E$ preserves adjoints. Taking the adjoint of
    $E(X)=\mu X$ gives the result.
\end{proof}

\begin{theorem}[Powers of a transfer-map eigenmatrix]
    \label{thm:transfer-map-eigenmatrix-power}
    \lean{MPSTensor.transferMap_pow_smul_eigenvector}
    \leanok
    \uses{def:transfer_map}
    If $\E_A(X)=\mu X$, then $\E_A^n(X)=\mu^nX$ for every $n\geq0$.
\end{theorem}

\begin{proof}\leanok
    \uses{thm:pow-on-eigenspace}
    This is Theorem~\ref{thm:pow-on-eigenspace} for the eigenspace of $\mu$.
\end{proof}

\begin{theorem}[Finite-order eigenmatrix is power-fixed]
    \label{thm:root-eigenvector-power-fixed}
    \lean{Kraus.pow_eigenvector_of_root}
    \leanok
    Let $E:\MN{D}\to\MN{D}$ be linear. If $E(X)=\mu X$ and $\mu^p=1$, then
    $E^p(X)=X$.
\end{theorem}

\begin{proof}\leanok
    \uses{thm:pow-on-eigenspace}
    Theorem~\ref{thm:pow-on-eigenspace} gives $E^p(X)=\mu^pX=X$.
\end{proof}

\begin{theorem}[Conjugate finite-order eigenmatrix is power-fixed]
    \label{thm:positive-map-conj-power-fixed}
    \lean{Kraus.pow_conjTranspose_eigenvector_of_root}
    \leanok
    \uses{def:positive_map}
    Let $E:\MN{D}\to\MN{D}$ be positive. If $E(X)=\mu X$ and $\mu^p=1$,
    then $E^p(X^\dagger)=X^\dagger$.
\end{theorem}

\begin{proof}\leanok
    \uses{thm:positive-map-conj-eigenvector,
        thm:root-eigenvector-power-fixed}
    The conjugate eigenvalue is $\overline\mu$, and
    $(\overline\mu)^p=\overline{\mu^p}=1$.
\end{proof}

\begin{theorem}[Trace of a non-fixed eigenmatrix]
    \label{thm:trace-preserving-eigenvector-trace-zero}
    \lean{Kraus.trace_eigenvector_eq_zero}
    \leanok
    \uses{def:trace_preserving}
    Let $E:\MN{D}\to\MN{D}$ preserve the trace. If $E(X)=\mu X$ and
    $\mu\neq1$, then $\tr(X)=0$.
\end{theorem}

\begin{proof}\leanok
    Trace preservation gives $\mu\tr(X)=\tr(E(X))=\tr(X)$. Since
    $\mu-1\neq0$, the trace vanishes.
\end{proof}

\begin{theorem}[Hermitian power-fixed point from a finite-order eigenmatrix]
    \label{thm:kraus-hermitian-power-fixed-point}
    \lean{Kraus.exists_hermitian_ne_zero_trace_zero_pow_fixedPoint}
    \leanok
    \uses{def:channel}
    Let $E$ be a channel, and suppose $E(X)=\mu X$ with $X\neq0$,
    $\mu\neq1$, and $\mu^p=1$. Then there is a nonzero Hermitian matrix $H$
    such that $\tr(H)=0$ and $E^p(H)=H$.
\end{theorem}

\begin{proof}\leanok
    \uses{thm:root-eigenvector-power-fixed,
        thm:positive-map-conj-power-fixed,
        thm:trace-preserving-eigenvector-trace-zero}
    At least one of $X+X^\dagger$ and $i(X^\dagger-X)$ is nonzero. Both are
    Hermitian, both have trace zero, and each is fixed by $E^p$.
\end{proof}

\begin{theorem}[Uniqueness of power-fixed positive matrices for a Kraus map]
    \label{thm:kraus-power-fixed-point-unique-from-spreading}
    \lean{Kraus.posSemidef_pow_fixedPoint_unique}
    \leanok
    \uses{def:kraus_map, def:vector_spread_span}
    Suppose $H_q(K,\varphi)=\C^D$ for every nonzero $\varphi$. If $p>0$ and
    $\rho,\sigma$ are nonzero positive-semidefinite fixed points of
    $\E_K^p$, then $\sigma=c\rho$ for some $c\in\C$.
\end{theorem}

\begin{proof}\leanok
    \uses{thm:kraus_power_fixed_point_posdef_from_spreading,
        thm:posdef_fixedpoint_proportional}
    Fixed-length spreading makes every nonzero positive-semidefinite fixed
    point of $\E_K^p$ positive definite. The critical-scalar argument then
    gives proportionality.
\end{proof}

\begin{theorem}[Powers of channels]
    \label{thm:channel-natural-power}
    \lean{Kraus.isChannel_pow}
    \leanok
    \uses{def:channel}
    If $E$ is a channel, then $E^p$ is a channel for every $p\geq0$.
\end{theorem}

\begin{proof}\leanok
    Complete positivity and trace preservation are both closed under
    composition and hold for the identity map.
\end{proof}

\begin{theorem}[Vanishing of Hermitian power-fixed points]
    \label{thm:kraus-hermitian-power-fixed-point-zero}
    \lean{Kraus.hermitian_pow_fixedPoint_eq_zero}
    \leanok
    \uses{def:kraus_map, def:vector_spread_span}
    Let $K$ be trace-preserving and suppose
    $H_q(K,\varphi)=\C^D$ for every nonzero $\varphi$. If $p>0$ and
    $H=H^\dagger$ satisfies $\tr(H)=0$ and $\E_K^p(H)=H$, then $H=0$.
\end{theorem}

\begin{proof}\leanok
    \uses{thm:channel-natural-power,
        thm:cesaro_fixedpoint,
        thm:kraus-power-fixed-point-unique-from-spreading,
        thm:positive_fixedpoint_decomp}
    Decompose $H$ as a difference of positive-semidefinite fixed points of the
    channel $\E_K^p$. Both are proportional to one nonzero fixed point. Their
    traces and $\tr(H)=0$ force equal coefficients, hence $H=0$.
\end{proof}

\begin{theorem}[Hermitian power-fixed point from a peripheral eigenvector]
    \label{thm:hermitian-power-fixed-point}
    \lean{MPSTensor.exists_hermitian_ne_zero_trace_zero_pow_fixedPoint}
    \leanok
    \uses{def:transfer_map}
    Let $A$ be a normalized MPS tensor, and suppose
    $\E_A(X)=\mu X$ with $X\neq0$, $\mu\neq1$, and $\mu^p=1$. Then there is a
    nonzero Hermitian matrix $H$ such that
    \begin{align}
        \tr(H)
         & =0,
        \notag\\
        \E_A^p(H)
         & =H.
        \notag
    \end{align}
    Moreover, $H$ is not positive semidefinite.
\end{theorem}

\begin{proof}\leanok
    \uses{thm:kraus-hermitian-power-fixed-point}
    Apply Theorem~\ref{thm:kraus-hermitian-power-fixed-point} to the transfer
    channel. The resulting matrix is nonzero, Hermitian, trace zero, and fixed
    by $\E_A^p$. A nonzero positive-semidefinite matrix has positive trace, so
    this matrix is not positive semidefinite.
\end{proof}

\begin{theorem}[Uniqueness of power-fixed positive matrices]
    \label{thm:mps-power-fixed-point-unique-from-spreading}
    \lean{MPSTensor.posSemidef_pow_fixedPoint_unique_of_isPrimitivePaper}
    \leanok
    \uses{def:transfer_map, def:vector_spread_span}
    Suppose there is a length $q\geq0$ such that
    $H_q(A,\varphi)=\C^D$ for every nonzero $\varphi\in\C^D$. Let $p>0$.
    If $\rho$ and $\sigma$ are nonzero positive-semidefinite fixed points of
    $\E_A^p$, then $\sigma=c\rho$ for some $c\in\C$.
\end{theorem}

\begin{proof}\leanok
    \uses{thm:kraus-power-fixed-point-unique-from-spreading}
    This is Theorem~\ref{thm:kraus-power-fixed-point-unique-from-spreading}
    for the Kraus family formed by the matrices of the tensor.
\end{proof}

\begin{theorem}[Positive-definite fixed point from vector spreading]
    \label{thm:kraus_posdef_fixed_point_from_tp_spreading}
    \lean{Kraus.exists_posDef_fixedPoint_of_isTP_of_vectorSpreadSpan_eq_top}
    \leanok
    \uses{def:kraus_map, def:vector_spread_span}
    Let $D>0$, and let $K=\{K^i\}_{i=0}^{d-1}$ be a finite family of
    $D\times D$ matrices satisfying
    \begin{align}
        \sum_{i=0}^{d-1}(K^i)^\dagger K^i
         & =\Id.
        \notag
    \end{align}
    Suppose there is a common length $q\ge0$ such that
    $H_q(K,\varphi)=\C^D$ for every nonzero $\varphi\in\C^D$.
    Then there is a positive-definite matrix $\rho$ such that
    $\E_K(\rho)=\rho$.
\end{theorem}

\begin{proof}\leanok
    \uses{thm:cesaro_fixedpoint,
        thm:kraus_fixed_point_posdef_from_spreading}
    The normalization makes $\E_K$ a channel. The Ces\`aro fixed-point theorem,
    Theorem~\ref{thm:cesaro_fixedpoint}, gives a nonzero positive-semidefinite
    fixed point. The common fixed-length spreading hypothesis makes this fixed
    point positive definite by
    Theorem~\ref{thm:kraus_fixed_point_posdef_from_spreading}.
\end{proof}

\begin{theorem}[Primitivity from vector spreading]
    \label{thm:kraus_primitive_from_tp_spreading}
    \lean{Kraus.isPrimitive_mapLM_of_isTP_of_vectorSpreadSpan_eq_top}
    \leanok
    \uses{def:kraus_map, def:vector_spread_span, def:primitive_channel}
    Let $D>0$, and let $K=\{K^i\}_{i=0}^{d-1}$ be a finite family of
    $D\times D$ matrices satisfying
    \begin{align}
        \sum_{i=0}^{d-1}(K^i)^\dagger K^i
         & =\Id.
        \notag
    \end{align}
    Suppose there is a common length $q\ge0$ such that
    $H_q(K,\varphi)=\C^D$ for every nonzero $\varphi\in\C^D$.
    Then the Kraus map $\E_K$ is primitive: it has a nonzero fixed point and
    $1$ is its only peripheral eigenvalue.
\end{theorem}

\begin{proof}\leanok
    \uses{thm:cesaro_fixedpoint, thm:kraus_irreducible_from_spreading,
        thm:peripheral_roots_irreducible_channel,
        thm:kraus-hermitian-power-fixed-point,
        thm:kraus-hermitian-power-fixed-point-zero}
    The normalization makes $\E_K$ a channel, so
    Theorem~\ref{thm:cesaro_fixedpoint} gives a nonzero positive-semidefinite
    fixed point. This supplies the peripheral eigenvalue $1$.
    Theorem~\ref{thm:kraus_irreducible_from_spreading} gives irreducibility.
    Thus $\E_K$ is an irreducible channel. By
    Theorem~\ref{thm:peripheral_roots_irreducible_channel}, every peripheral
    eigenvalue $\mu$ is a root of unity; choose $p>0$ with
    $\mu^p=1$.

    Suppose $\mu\neq1$. Theorem~\ref{thm:kraus-hermitian-power-fixed-point}
    gives a nonzero Hermitian matrix $H$ with $\tr(H)=0$ and
    $\E_K^p(H)=H$. Theorem~\ref{thm:kraus-hermitian-power-fixed-point-zero}
    gives $H=0$, a contradiction. Thus $\mu=1$.
\end{proof}

\begin{theorem}[Strengthened form of the primitivity equivalence]
    \label{thm:wielandt_proposition_three}
    \lean{MPSTensor.primitivePaper_iff_hasEventuallyFullKrausRank}
    \lean{MPSTensor.primitivePaper_iff_stronglyIrreducible}
    \lean{MPSTensor.hasEventuallyFullKrausRank_iff_stronglyIrreducible}
    \leanok
    \uses{def:paper_primitive, def:eventually_full_kraus_rank,
        def:strongly_irreducible_paper}
    Let $D>0$ and suppose
    $\sum_i(A^i)^\dagger A^i=\Id$. The following are equivalent:
    \begin{enumerate}
        \item there is a positive $q$ for which
              $H_q(A,\varphi)=\C^D$ for every $\varphi\neq0$;
        \item there is a positive $i$ for which $S_i(A)=\MN{D}$;
        \item $\E_A$ is strongly irreducible.
    \end{enumerate}
    Proposition~3 of~\cite{Sanz2010Quantum} states the same equivalence with
    item~(3) replaced by the literal source condition: $1$ is the only
    peripheral eigenvalue, and its one-dimensional eigenspace is generated by
    a positive-definite fixed point. Here item~(3) includes the additional
    irreducibility condition.
\end{theorem}

\begin{proof}\leanok
    \uses{thm:blocking_primitivity}
    If $S_i(A)=\MN{D}$ and $\varphi\neq0$, then
    $S_i(A)\varphi=\C^D$, proving (2)$\Rightarrow$(1).
    For (1)$\Rightarrow$(3), first obtain a nonzero positive fixed point
    $\rho$. The spreading condition makes $\rho$ positive definite and makes
    the tensor irreducible. The blocking-periodicity theorem then gives
    $p>0$ such that $\E_A^p$ is primitive. If
    $\E_A(X)=\lambda X$ with $|\lambda|=1$, then
    $\E_A^p(X)=\lambda^pX$; uniqueness of the peripheral eigenvalue of the
    primitive power gives $\lambda^p=1$. If $\lambda\neq1$, the eigenvector
    yields a nonzero trace-zero Hermitian matrix $H$ fixed by $\E_A^p$.
    Decompose $H=Q_1-Q_2$ into positive semidefinite $\E_A^p$-fixed parts.
    Paper primitivity makes every nonzero positive fixed point of the power
    proportional to $\rho$, while $\tr(H)=0$ makes the two proportionality
    constants equal. Hence $H=0$, a contradiction. Thus every peripheral
    eigenvalue equals $1$, and $\E_A$ is strongly irreducible.

    For (3)$\Rightarrow$(2), write
    $\E_A^m=P_\rho+(\E_A-P_\rho)^m$. The complementary powers tend to zero.
    If $S_m(A)\neq\MN{D}$ for every $m$, choose a nonzero matrix $B$
    orthogonal to $S_m(A)$. The associated trace pairing vanishes on the
    left, while positive definiteness gives a lower bound
    $c\lVert B\rVert^2$ for the $P_\rho$ term and the complementary term is
    $o(\lVert B\rVert^2)$. For large $m$ this is a contradiction.
\end{proof}

% Scope restriction (primitive channels): Proposition 1 is universal, whereas
% the theorem below assumes paper primitivity. See
% docs/paper-gaps/quantum_wielandt_deviation_v1.tex.
\begin{theorem}[Primitive-channel specialization of the index bound]
    \label{thm:q_index_le_i_index}
    \lean{MPSTensor.qIndex_le_iIndex_of_isPrimitivePaper}
    \leanok
    \uses{def:wielandt_i_index, def:wielandt_q_index,
        thm:wielandt_proposition_three}
    Under the hypotheses of
    Theorem~\ref{thm:wielandt_proposition_three}, paper primitivity implies
    \begin{align}
        q(\E_A) & \le\iota(A).
        \notag
    \end{align}
    Proposition~1 of~\cite{Sanz2010Quantum} states this inequality for every
    quantum channel; the present theorem records its specialization under
    normalization and paper primitivity.
\end{theorem}

\begin{proof}\leanok
    By Theorem~\ref{thm:wielandt_proposition_three}, $\iota(A)$ is a
    positive length with $S_{\iota(A)}(A)=\MN{D}$. Therefore
    $H_{\iota(A)}(A,\varphi)=\C^D$ for every $\varphi\neq0$, so the defining
    minimum for $q(\E_A)$ is at most $\iota(A)$.
\end{proof}

\begin{theorem}[Unique fixed points of full rank]
    \label{thm:wolf_6_15}
    \lean{MPSTensor.wolf_theorem_6_15}
    \leanok
    \uses{thm:wielandt_proposition_three}
    Let $D>0$ and suppose $\sum_i(A^i)^\dagger A^i=\Id$. If there is a
    positive $n$ for which $S_n(A)=\MN{D}$, then there is a unique positive
    definite density matrix $\rho\in\MN{D}$ such that every fixed point of
    $\E_A$ is a scalar multiple of $\rho$.

    This is~\cite[Theorem~6.15]{Wolf2012Quantum}.
\end{theorem}

\begin{proof}\leanok
    \uses{thm:wielandt_proposition_three, lem:fixedpoint_unique_spectral}
    By Theorem~\ref{thm:wielandt_proposition_three}, the span hypothesis
    makes $\E_A$ strongly irreducible, giving a positive-definite fixed
    point $\rho_0$ together with peripheral primitivity (unique peripheral
    eigenvalue $1$) and irreducibility of $\E_A$. Peripheral primitivity and
    irreducibility of $\E_A$ together give the complementary transfer-map
    gap at $\rho_0$ (Lemma~\ref{lem:fixedpoint_unique_spectral}): the
    restriction of $\E_A$ off the fixed-point
    projection at $\rho_0$ has spectral radius strictly less than one. This
    gap forces every fixed point of $\E_A$ to be a scalar multiple of
    $\rho_0$; rescaling $\rho_0$ by the inverse of its (positive) trace
    gives the unique density-matrix representative $\rho$.

    The source proves this instead by passing to $\E_A^n$, showing every
    $\E_A^n(\ket\psi\!\bra\psi)$ has full rank from the span hypothesis, and
    ruling out a fixed-point space of dimension larger than one via
    Corollary~6.5 of~\cite{Wolf2012Quantum}. Both routes use only Wolf's own
    Chapter~6 machinery; the present proof reuses the already-formalized
    Proposition~3 equivalence.
\end{proof}

\begin{lemma}[Monotonicity of cumulative span]\label{lem:cumulative_mono}
    \lean{MPSTensor.cumulativeSpan_mono}
    \leanok
    \uses{def:cumulative_span}
    For every $n$, one has $T_n(A)\le T_{n+1}(A)$.
\end{lemma}

\begin{proof}\leanok\uses{def:cumulative_span}
    The span $T_{n+1}$ includes all words of length at most $n+1$ and
    therefore contains all words of length at most~$n$.
\end{proof}

\begin{lemma}[Dimension bound]\label{lem:cumulative_finrank_le}
    \lean{MPSTensor.cumulativeSpan_finrank_le}
    \leanok
    \uses{def:cumulative_span}
    For every $n$, one has $\dim T_n(A)\le D^2$.
\end{lemma}

\begin{proof}\leanok
    Since $T_n(A)\subseteq\MN{D}$ and $\dim\MN{D}=D^2$, the result follows.
\end{proof}

\begin{theorem}[Cumulative span stabilisation]\label{thm:cumulative_stable}
    \lean{MPSTensor.cumulativeSpan_stable}
    \leanok
    \uses{def:cumulative_span}
    If $T_n(A)=T_{n+1}(A)$, then $T_m(A)=T_n(A)$ for every $m\ge n$.
\end{theorem}

\begin{proof}
    \leanok
    If $T_n=T_{n+1}$, then $S_{n+1}\subseteq T_n$.
    Left-multiplying any element of $S_{n+1}$ by $A^i$ gives
    an element of $S_{n+2}$, so $S_{n+2}\subseteq T_{n+1}=T_n$.
    By induction, $S_m\subseteq T_n$ for every $m>n$, hence $T_m=T_n$.
\end{proof}

\begin{theorem}[Dimension growth]\label{thm:cumulative_strict_mono}
    \lean{MPSTensor.cumulativeSpan_finrank_strict_mono}
    \leanok
    \uses{def:cumulative_span}
    If $T_n(A)\neq T_{n+1}(A)$, then
    $\dim T_{n+1}(A)>\dim T_n(A)$.
\end{theorem}

\begin{proof}
    \leanok
    \uses{lem:cumulative_mono}
    The inclusion $T_n\le T_{n+1}$ from
    Lemma~\ref{lem:cumulative_mono}, together with
    $T_n\neq T_{n+1}$, gives strict subspace inclusion and hence strict
    dimension growth.
\end{proof}

\begin{lemma}[Word span characterises block injectivity]\label{lem:wordspan_blk_inj}
    \lean{MPSTensor.wordSpan_eq_top_iff_isNBlkInjective}
    \leanok
    \uses{def:word_span, def:l_blk_injective}
    One has $S_L(A)=\MN{D}$ if and only if $A$ is $L$-block injective.
\end{lemma}

\begin{proof}\leanok\uses{def:word_span, def:l_blk_injective}
    Both conditions assert that the span of all products of length~$L$
    equals $\MN{D}$.
\end{proof}

\begin{theorem}[Positive-length propagation of block injectivity]
    \label{thm:wordspan_blk_inj_succ}
    \lean{MPSTensor.wordSpan_succ_eq_mul_left}
    \lean{MPSTensor.isNBlkInjective_succ_of_isNBlkInjective}
    \lean{MPSTensor.isNBlkInjective_of_le}
    \leanok
    \uses{def:word_span, def:l_blk_injective}
    For every $n$,
    \begin{align}
        S_{n+1}(A)
         & =\spn_{\C}\{A^i\}_{i=0}^{d-1}S_n(A).
        \notag
    \end{align}
    Consequently, if $L>0$ and $A$ is $L$-block injective, then
    $A$ is $m$-block injective for every $m\ge L$.
    This auxiliary length-shift statement combines the block-injectivity
    discussion in~\cite[Section~II]{Cirac2016MPDO_arXiv} with
    \cite[Appendix~C.3, Lemma~L]{Cirac2016MPDO_arXiv}; it is not stated there
    as a separate lemma.
\end{theorem}

\begin{proof}\leanok\uses{def:word_span}
    Every word of length $n+1$ factors into its first letter and a word
    of length~$n$, which proves the displayed equality by bilinearity.
    Write $L=n+1$. Since $S_L(A)=\MN{D}$, one has
    $S_n(A)\le S_L(A)$. Hence
    \begin{align}
        \MN{D}
        =S_L(A)
        =\spn_{\C}\{A^i\}_{i=0}^{d-1}S_n(A)
        \le\spn_{\C}\{A^i\}_{i=0}^{d-1}S_L(A)
        =S_{L+1}(A).
        \notag
    \end{align}
    Iterating this implication gives $S_m(A)=\MN{D}$ for every $m\ge L$.
\end{proof}

\section{Nonzero trace product}

\begin{theorem}[Cumulative span reaches $\MN{D}$]\label{thm:cumulative_eq_top}
    \label{thm:wielandt_bound}
    \lean{MPSTensor.cumulativeSpan_eq_top}
    \leanok
    \uses{def:cumulative_span, def:normal}
    Let $D>0$. If $A$ is normal, then $T_{D^2}(A)=\MN{D}$.
    This cumulative estimate supports the proof of
    \cite[Lemma~1]{Sanz2010Quantum}; it is not a conclusion of source
    Theorem~1.
\end{theorem}

\begin{proof}\leanok\uses{thm:cumulative_eq_top_bound}
    This is Theorem~\ref{thm:cumulative_eq_top_bound}.
\end{proof}

\begin{theorem}[Cumulative span reaches $\MN{D}$ --- sharp bound]%
    \label{thm:cumulative_span_sharp}
    \lean{MPSTensor.cumulativeSpan_eq_top_of_isNormal_sharp}
    \leanok
    \uses{def:cumulative_span, def:normal}
    Let $D>0$. If $A$ is normal, then
    $T_{D^2-d'+1}(A)=\MN{D}$, where
    $d'=\dim S_1(A)=\krausrk(A)$.
    This is the cumulative-span estimate used to prove the sharp trace
    conclusion in \cite[Lemma~1]{Sanz2010Quantum}.
\end{theorem}

\begin{proof}
    \leanok
    \uses{thm:cumulative_strict_mono, lem:cumulative_finrank_le}
    If $T_n(A)\neq T_{n+1}(A)$, then the dimension strictly increases.
    Since $S_1(A)\subseteq T_1(A)$ and $\dim S_1(A)=d'$, one has
    $\dim T_1(A)\ge d'$.
    After at most $D^2-d'$ additional strict-growth steps, the dimension
    reaches $D^2=\dim\MN{D}$.
    Stabilisation before reaching $\MN{D}$ contradicts normality.
    Hence $T_{D^2-d'+1}(A)=\MN{D}$.
\end{proof}

\begin{theorem}[Eigenvector extraction from a full cumulative span]
    \label{thm:eigenvector_from_full_cumulative_span}
    \lean{Kraus.exists_eigenvector_of_cumulativeSpan_eq_top}
    \lean{MPSTensor.exists_eigenvector_of_cumulativeSpan_eq_top}
    \leanok
    \uses{def:cumulative_span}
    Let $D>0$ and suppose $T_N(K)=\MN{D}$. Then there are a word $w$, a
    nonzero scalar $\mu$, and a nonzero vector $\varphi\in\C^D$ such that
    \begin{align}
        |w| & \le N, \notag\\
        K^w\varphi & = \mu\varphi. \notag
    \end{align}
\end{theorem}

\begin{proof}
    \leanok
    \uses{thm:eigenvector_from_trace}
    Since the identity lies in $T_N(K)$ and has nonzero trace, some word of
    length at most $N$ has nonzero trace. Such a matrix has a nonzero
    eigenvalue and a corresponding nonzero eigenvector by
    Theorem~\ref{thm:eigenvector_from_trace}.
\end{proof}

\begin{theorem}[Sharp nonzero trace product for a primitive channel]%
    \label{thm:nonzero_trace_word_primitive_sharp}
    \lean{MPSTensor.exists_nonzero_trace_word_of_isPrimitivePaper_sharp_pos}
    \leanok
    \uses{def:paper_primitive, def:wielandt_kraus_rank,
        thm:wielandt_proposition_three}
    Let $D>0$, suppose $\sum_i(A^i)^\dagger A^i=\Id$, and assume that
    $A$ is primitive in the sense of
    Definition~\ref{def:paper_primitive}. Then there is a word $w$ with
    \begin{align}
        1\le|w|  & \le D^2-\krausrk(A)+1,
        \notag                            \\
        \tr(A^w) & \neq0.
        \notag
    \end{align}
    This is the sharp trace conclusion of
    \cite[Lemma~1]{Sanz2010Quantum}.
\end{theorem}

\begin{proof}\leanok
    Theorem~\ref{thm:wielandt_proposition_three} gives eventual full Kraus
    rank, hence normality. Put $b=D^2-\krausrk(A)+1$ and let
    $V_n$ be the span of $A^w$ with $1\le|w|\le n$. Then
    $S_1(A)\subseteq V_1$, so $\dim V_1\ge\krausrk(A)$. If
    $V_n=V_{n+1}$ for some $n\ge1$, then $V_n$ is closed under left
    multiplication by every $A^i$. It therefore contains every positive-length
    word product; normality forces $V_n=\MN{D}$. Thus, until the full matrix
    algebra is reached, the dimensions increase strictly, and the bound on
    $\dim\MN{D}$ gives $V_b=\MN{D}$. Express $\Id$ as a linear combination of
    the positive-length words defining $V_b$. Since $\tr(\Id)=D\neq0$, at
    least one has nonzero trace.
\end{proof}

\begin{theorem}[Sharp nonzero trace product]%
    \label{thm:nonzero_trace_word_sharp}
    \lean{MPSTensor.exists_nonzero_trace_word_sharp}
    \leanok
    \uses{def:normal}
    Let $D>0$. If $A$ is normal, then there exists a word $w$ of length at
    most $D^2-\krausrk(A)+1$ such that $\tr(A^w)\neq0$.
    This normality variant has the same sharp length bound as
    \cite[Lemma~1]{Sanz2010Quantum}; the source hypotheses are stated in
    Theorem~\ref{thm:nonzero_trace_word_primitive_sharp}.
\end{theorem}

\begin{proof}
    \leanok
    \uses{thm:cumulative_span_sharp}
    By Theorem~\ref{thm:cumulative_span_sharp}, the identity $\Id$ lies in
    $T_{D^2-d'+1}(A)$ and is therefore a linear combination of word
    products of length at most $D^2-d'+1$.
    At least one such product has nonzero trace.
\end{proof}

\begin{theorem}[Sharp nonzero trace product --- positive length]%
    \label{thm:nonzero_trace_word_sharp_pos}
    \lean{MPSTensor.exists_nonzero_trace_word_sharp_pos}
    \leanok
    \uses{def:normal, thm:nonzero_trace_word_sharp}
    If $A$ is normal and $D\ge1$, then there exists a word~$w$ of
    \emph{positive} length
    $1\le|w|\le D^2-\krausrk(A)+1$ such that $\tr(A^w)\neq0$.
    This strengthens Theorem~\ref{thm:nonzero_trace_word_sharp} by requiring
    $|w|\ge1$, a feature needed for the blocking argument in the general
    Wielandt bound below.
\end{theorem}

\begin{proof}
    \leanok
    \uses{thm:cumulative_span_sharp}
    Define the \emph{positive-level} cumulative span
    $V_n=\spn_{\C}\{A^w:1\le|w|\le n\}$.
    The same stabilization-or-strict-growth dichotomy shows
    $V_{D^2-d'+1}=\MN{D}$: if $V_n$ stabilises and is not $\MN{D}$, then
    $V_n$ is closed under left multiplication by each $A^i$, so
    $\wordspn{N}{A}\subseteq V_n\neq\MN{D}$ for every $N\ge1$,
    contradicting normality.
    Since $V_{D^2-d'+1}=\MN{D}$, the identity $\Id$ lies in
    $V_{D^2-d'+1}$ and is therefore a linear combination of word products
    of positive length. At least one such product has nonzero trace because
    $\tr(\Id)=D\neq0$ for $D\ge1$.
\end{proof}

\section{Eigenvector spreading}

\begin{definition}[Cumulative vector span]\label{def:cumulative_vector_span}
    \lean{MPSTensor.cumulativeVectorSpan}
    \leanok
    \uses{def:mps_tensor}
    Given an MPS tensor~$A$ and a vector $\varphi\in\C^D$, the
    \emph{cumulative vector span} is
    \begin{align}
        K_n(A,\varphi)
         & =\spn_{\C}\{A^w\varphi:|w|\le n\}\subseteq\C^D.
        \notag
    \end{align}
    This is~\cite[proof of Lemma~2(a)]{Sanz2010Quantum}.
\end{definition}

\begin{theorem}[Eigenvector spreading]\label{thm:eigenvector_spreading}
    \lean{MPSTensor.eigenvector_spreading}
    \leanok
    \uses{def:cumulative_vector_span, def:normal}
    Let $A$ be a normal MPS tensor and let
    $A^{i_0}\varphi=\mu\varphi$, where $\mu\neq0$ and $\varphi\neq0$.
    Then $K_{D-1}(A,\varphi)=\C^D$.
    This cumulative form supplies the dimension-growth step in
    \cite[Lemma~2(a)]{Sanz2010Quantum}.
\end{theorem}

\begin{proof}
    \leanok
    \uses{thm:cumulative_eq_top, lem:vector_span_from_matrix,
        thm:vector_span_stable}
    The vector span $K_n$ is monotone and $\dim K_n\le D$.
    If $K_n=K_{n+1}$ for some $n<D-1$, the span stabilises by
    Theorem~\ref{thm:vector_span_stable}.
    But $T_{D^2}(A)=\MN{D}$ by Theorem~\ref{thm:cumulative_eq_top}, so
    Lemma~\ref{lem:vector_span_from_matrix} gives
    $K_{D^2}(A,\varphi)=\C^D$, contradicting early stabilisation.
    Hence $\dim K_n$ grows by at least~$1$ at each step, reaching~$D$ by
    step $D-1$.
\end{proof}

\begin{theorem}[Fixed-length spreading from eventual word-span fullness]
    \label{thm:eigenvector_spreading_eventually_full}
    \lean{Kraus.vectorSpreadSpan_eq_top_of_hasEventuallyFullWordSpan_of_eigenvector}
    \leanok
    \uses{def:kraus_eventually_full_word_span, def:vector_spread_span}
    Let $K=(K^i)_{i=0}^{d-1}$ be a finite family of $D\times D$ matrices, where
    $D>0$, and suppose that $S_n(K)=\MN{D}$ for every sufficiently large
    $n$. If $\varphi\neq0$, $\mu\neq0$, and
    $K^{i_0}\varphi=\mu\varphi$, then
    $H_{D-1}(K,\varphi)=\C^D$.
\end{theorem}

\begin{proof}\leanok
    \uses{thm:eigenvector_spreading, lem:eigenvector_padding}
    Eventual fullness gives a level at which the cumulative matrix span is
    all of $\MN{D}$. The cumulative vector spaces therefore reach $\C^D$ by
    level $D-1$. The eigenvector relation then pads shorter words to length
    $D-1$ by nonzero powers of $\mu$.
\end{proof}

\begin{theorem}[Fixed-length eigenvector spreading for a primitive channel]
    \label{thm:eigenvector_spreading_primitive}
    \lean{MPSTensor.vectorSpreadSpan_eq_top_of_isPrimitivePaper_of_eigenvector}
    \leanok
    \uses{def:paper_primitive, def:vector_spread_span,
        thm:wielandt_proposition_three, thm:eigenvector_spreading,
        lem:eigenvector_padding}
    Let $D>0$, suppose $\sum_i(A^i)^\dagger A^i=\Id$, and assume that
    $A$ is primitive in the sense of
    Definition~\ref{def:paper_primitive}. If $\varphi\neq0$, $\mu\neq0$, and
    $A^{i_0}\varphi=\mu\varphi$, then
    $H_{D-1}(A,\varphi)=\C^D$.
    This is~\cite[Lemma~2(a)]{Sanz2010Quantum}.
\end{theorem}

\begin{proof}\leanok
    Proposition~3, in the form of
    Theorem~\ref{thm:wielandt_proposition_three}, gives eventual full Kraus
    rank and hence normality. The cumulative spaces $K_n(A,\varphi)$ cannot
    stabilise below $\C^D$, because every matrix eventually lies in an exact
    word span. Their dimensions therefore reach $D$ by $n=D-1$, giving
    $K_{D-1}(A,\varphi)=\C^D$. Finally, the eigenvector relation pads every
    shorter word to length $D-1$ by a nonzero power of $\mu$, so
    $K_{D-1}(A,\varphi)=H_{D-1}(A,\varphi)$.
\end{proof}

\section{Quantum Wielandt bounds}

Section~\ref{sec:app_wielandt_augmentation} proves the one-step augmentation
facts used in cases~(2) and~(3). The blocking identities used in the general
case are proved in Section~\ref{sec:app_wielandt_blocking_support}.

\begin{remark}[Explicit blocking-length estimates]\leanok
    Theorem~\ref{thm:wielandt_bound} gives the $D^2$ cumulative-span bound.
    The full-rank index bound
    $\iota(A)\le(D^2-d'+1)D^2$
    (\cite[Theorem~1, case~(1)]{Sanz2010Quantum}) is proved in
    Theorem~\ref{thm:wielandt_general} below.
    The left-canonical normal blocked-injectivity estimate $D^4$ is proved
    below in
    Theorem~\ref{thm:bounded_injective_blocking_normal_left_canonical}.
    A sharper $3D^5$ canonical-form block-separation estimate, developed in
    the full blueprint, combines this $D^4$ estimate with the direct-sum
    argument of \cite{PerezGarcia2007Matrix}.
\end{remark}

\begin{theorem}[Wielandt case (2), one-step invertible element]%
    \label{thm:wielandt_case2_one_step_span}
    \lean{MPSTensor.wordSpan_eq_top_of_isPrimitivePaper_of_mem_wordSpan_one_of_isUnit}
    \leanok
    \uses{def:word_span, def:one_step_augment,
        lem:one_step_augment_word_span, lem:one_step_augment_normal_kraus_rank,
        thm:wielandt_case2_normal_generator, def:paper_primitive}
    Let $A$ be an MPS tensor with bond dimension $D\ge1$, satisfying
    $\sum_i(A^i)^\dagger A^i=\Id$, and suppose that $A$ is primitive in the
    sense of Definition~\ref{def:paper_primitive}.
    If $S_1(A)$ contains an invertible matrix~$X$, then
    \begin{align}
        S_{D^2-\krausrk(A)+1}(A)
         & =\MN{D}.
        \label{eq:wld_case2_one_step_wordspan_top}
    \end{align}
    In the terminology of~\cite{Sanz2010Quantum}, this is the case where
    the first application space $S_1(A)$ contains an invertible operator.
    This is the exact word-span form of
    \cite[Theorem~1, case~(2)]{Sanz2010Quantum}.
\end{theorem}

\begin{proof}
    \leanok
    \uses{thm:wielandt_proposition_three,
        lem:one_step_augment_word_span, lem:one_step_augment_normal_kraus_rank,
        thm:wielandt_case2_normal_generator}
    Adjoin $X$ as an extra first Kraus operator, obtaining $\widetilde A$.
    Theorem~\ref{thm:wielandt_proposition_three} gives normality of $A$, hence normality of
    $\widetilde A$; the Kraus rank and all exact word spans are unchanged.
    The first Kraus operator of $\widetilde A$ is invertible, so
    Theorem~\ref{thm:wielandt_case2_normal_generator} gives
    $S_{D^2-\krausrk(\widetilde A)+1}(\widetilde A)=\MN{D}$.
    Transfer this equality back to $A$.
\end{proof}

\begin{corollary}[Wielandt case (2), index bound]%
    \label{cor:wielandt_case2_one_step_index}
    \lean{MPSTensor.iIndex_le_of_mem_wordSpan_one_of_isUnit}
    \leanok
    \uses{thm:wielandt_case2_one_step_span}
    Under the hypotheses of Theorem~\ref{thm:wielandt_case2_one_step_span},
    one has $\iota(A)\le D^2-\krausrk(A)+1$, where
    $\iota(A)=\min\{n\ge1:S_n(A)=\MN{D}\}$.
\end{corollary}

\begin{proof}
    \leanok
    \uses{thm:wielandt_case2_one_step_span}
    The word-span equality (\ref{eq:wld_case2_one_step_wordspan_top}) gives an
    admissible value in the defining minimum for $\iota(A)$.
\end{proof}

\begin{lemma}[Rank-one extraction from a non-invertible eigenvector]%
    \label{lem:rank_one_noninvertible_eigenvector}
    \lean{MPSTensor.vecMulVec_eigenvector_exact_wordSpan}
    \leanok
    \uses{def:word_span, def:normal}
    Let $A$ be a normal MPS tensor with bond dimension $D\ge1$.
    Suppose that a Kraus operator $A^{i_0}$ is non-invertible and that
    there exist a nonzero vector $\varphi\in\C^D$ and a nonzero scalar $\mu$
    such that $A^{i_0}\varphi=\mu\varphi$.
    Then, for every $\psi\in\C^D$,
    $\ketbra{\varphi}{\psi}\in S_{D^2-D+1}(A)$.
\end{lemma}

\begin{proof}
    \leanok
    Let $\mathcal L(M)=A^{i_0}M$, and choose $r$ so that
    $\mathcal L^r$ annihilates the generalized zero-eigenspace of $\mathcal L$;
    equivalently, $\range(\mathcal L^r)$ is its stabilized nonzero-spectral
    part. The dimension-growth argument gives a level $n_0$ for which
    \begin{align}
        (A^{i_0})^rS_{n_0}(A)
         & =\{(A^{i_0})^rM:M\in\MN{D}\}.
        \notag
    \end{align}
    Since $A^{i_0}\varphi=\mu\varphi$ and $\mu\neq0$, iterating gives
    $(A^{i_0})^r\varphi=\mu^r\varphi$ and hence
    $\varphi=\mu^{-r}(A^{i_0})^r\varphi\in\range((A^{i_0})^r)$.
    Therefore, for every $\psi\in\C^D$,
    $\ketbra{\varphi}{\psi}\in S_{r+n_0}(A)$.
    The corresponding Fitting-index estimate gives
    $r+n_0\le D^2-D+1$.
    Since
    \begin{align}
        A^{i_0}\ketbra{\varphi}{\psi}
         & =\ketbra{A^{i_0}\varphi}{\psi}
        =\mu\ketbra{\varphi}{\psi},
        \notag
    \end{align}
    one has
    $\ketbra{\varphi}{\psi}
        =\mu^{-1}A^{i_0}\ketbra{\varphi}{\psi}$.
    Thus
    $\ketbra{\varphi}{\psi}\in S_k(A)$ implies
    $\ketbra{\varphi}{\psi}\in S_{k+1}(A)$, because left multiplication
    by $A^{i_0}$ maps $S_k(A)$ into $S_{k+1}(A)$.
    Iterating from $k=r+n_0$ and using $r+n_0\le D^2-D+1$ gives
    $\ketbra{\varphi}{\psi}\in S_{D^2-D+1}(A)$.
\end{proof}

\begin{theorem}[Sharp rank-one extraction for a primitive channel]%
    \label{thm:rank_one_noninvertible_eigenvector_primitive}
    \lean{MPSTensor.vecMulVec_mem_wordSpan_of_isPrimitivePaper_of_noninvertible_eigenvector}
    \leanok
    \uses{def:paper_primitive, thm:wielandt_proposition_three,
        def:fitting_decomp, thm:nilpotent_pow_bound}
    Let $D>0$, suppose $\sum_i(A^i)^\dagger A^i=\Id$, and assume that
    $A$ is primitive in the sense of
    Definition~\ref{def:paper_primitive}. Suppose $A^{i_0}$ is not invertible
    and
    $A^{i_0}\varphi=\mu\varphi$ with $\varphi\neq0$ and $\mu\neq0$.
    Then, for every $\psi\in\C^D$,
    \begin{align}
        \ketbra{\varphi}{\psi}
         & \in S_{D^2-D+1}(A).
        \notag
    \end{align}
    This is the quantitative rank-one conclusion of
    \cite[Lemma~2(b)]{Sanz2010Quantum}.
\end{theorem}

\begin{proof}\leanok
    Theorem~\ref{thm:wielandt_proposition_three} first gives normality.
    Apply the Fitting decomposition to left multiplication
    $\mathcal L(M)=A^{i_0}M$. Choose $r$ so that $\mathcal L^r$ kills its
    generalized zero-eigenspace. Dimension growth then gives $n_0$ with
    \begin{align}
        (A^{i_0})^rS_{n_0}(A)
              & =\{(A^{i_0})^rM:M\in\MN{D}\},
        \notag                                \\
        r+n_0 & \le D^2-D+1.
        \notag
    \end{align}
    Since $(A^{i_0})^r\varphi=\mu^r\varphi$ and $\mu\neq0$, the vector
    $\varphi$ lies in the range of $(A^{i_0})^r$. Hence
    $\ketbra{\varphi}{\psi}\in S_{r+n_0}(A)$ for every $\psi$.
    The identity
    $\ketbra{\varphi}{\psi}=\mu^{-1}A^{i_0}
        \ketbra{\varphi}{\psi}$
    pads this element to the exact length $D^2-D+1$.
\end{proof}

\begin{theorem}[Wielandt case (3), one-step non-invertible element]%
    \label{thm:wielandt_case3_one_step_span}
    \lean{MPSTensor.wordSpan_eq_top_of_isPrimitivePaper_of_mem_wordSpan_one_of_noninvertible_eigenvector}
    \leanok
    \uses{def:word_span, def:one_step_augment,
        lem:one_step_augment_word_span, lem:one_step_augment_normal_kraus_rank,
        thm:eigenvector_spreading, lem:rank_one_noninvertible_eigenvector,
        thm:lemma2b_assembly, def:paper_primitive}
    Let $A$ be an MPS tensor with bond dimension $D\ge1$, satisfying
    $\sum_i(A^i)^\dagger A^i=\Id$, and suppose that $A$ is primitive in the
    sense of Definition~\ref{def:paper_primitive}.
    If $S_1(A)$ contains a non-invertible matrix $X$ with a nonzero
    eigenvalue, then $S_{D^2}(A)=\MN{D}$.
    More explicitly, it is enough to have $X\varphi=\mu\varphi$ with
    $\varphi\neq0$ and $\mu\neq0$.
    In the terminology of~\cite{Sanz2010Quantum}, this is the case where
    the first application space $S_1(A)$ contains a non-invertible operator
    with a nonzero eigenvalue.
    This is the exact word-span form of
    \cite[Theorem~1, case~(3)]{Sanz2010Quantum}.
\end{theorem}

\begin{proof}
    \leanok
    \uses{thm:wielandt_proposition_three,
        lem:one_step_augment_word_span, lem:one_step_augment_normal_kraus_rank,
        thm:eigenvector_spreading, lem:rank_one_noninvertible_eigenvector,
        thm:lemma2b_assembly}
    Theorem~\ref{thm:wielandt_proposition_three} gives normality.
    Adjoin $X$ as an extra first Kraus operator. The augmented tensor remains
    normal and has the same exact word spans as $A$.
    The eigenvector relation for $X$ is now a one-generator eigenvector
    relation. The eigenvector-spreading argument gives full fixed-length
    vector span by length $D-1$, and
    Lemma~\ref{lem:rank_one_noninvertible_eigenvector} gives the rank-one
    operators $\ketbra{\varphi}{\psi}$ at length $D^2-D+1$.
    Multiplying these two fixed-length spans gives $S_{D^2}$ for the
    augmented tensor, hence for $A$.
\end{proof}

\begin{corollary}[Wielandt case (3), index bound]%
    \label{cor:wielandt_case3_one_step_index}
    \lean{MPSTensor.iIndex_le_sq_of_mem_wordSpan_one_of_noninvertible_eigenvector}
    \leanok
    \uses{thm:wielandt_case3_one_step_span}
    Under the hypotheses of Theorem~\ref{thm:wielandt_case3_one_step_span},
    one has $\iota(A)\le D^2$.
\end{corollary}

\begin{proof}
    \leanok
    \uses{thm:wielandt_case3_one_step_span}
    The equality $S_{D^2}(A)=\MN{D}$ gives an admissible value in the
    defining minimum for $\iota(A)$.
\end{proof}

\begin{theorem}[Wielandt's inequality --- general bound]%
    \label{thm:wielandt_general}
    \lean{MPSTensor.iIndex_le_general_of_isPrimitivePaper}
    \leanok
    \uses{def:paper_primitive, def:wielandt_i_index,
        def:wielandt_kraus_rank}
    Let $A$ be a normalized MPS tensor with $D>0$,
    $\sum_i(A^i)^\dagger A^i=\Id$, that is primitive in the sense of
    Definition~\ref{def:paper_primitive}.
    Then
    \begin{align}
        \iota(A)
         & \le(D^2-\krausrk(A)+1)D^2,
        \notag
    \end{align}
    where $\iota(A)=\min\{n\ge1:S_n(A)=\MN{D}\}$ is the
    full-Kraus-rank index.
    This is~\cite[Theorem~1, case~(1)]{Sanz2010Quantum}.
\end{theorem}

\begin{proof}
    \leanok
    \uses{thm:wielandt_proposition_three,
        thm:nonzero_trace_word_primitive_sharp, thm:eigenvector_from_trace,
        thm:isNormal_blockTensor, thm:wielandt_case2_normal_generator,
        cor:wielandt_case3_generator}
    Theorem~\ref{thm:wielandt_proposition_three} gives normality.
    By Theorem~\ref{thm:nonzero_trace_word_primitive_sharp}, there exists a
    \emph{positive-length} word~$w$ with
    $|w|\le D^2-d'+1$ and $\tr(A^w)\neq0$, where
    $d'=\krausrk(A)$.
    The matrix $A^w$ has nonzero trace, hence a nonzero eigenvalue~$\mu$
    and eigenvector~$\varphi$ by
    Theorem~\ref{thm:eigenvector_from_trace}.
    Set $n=|w|$ and let $B=A^{[n]}$ be the blocked tensor.
    Then $B$ is normal by Theorem~\ref{thm:isNormal_blockTensor}, and the
    encoding of $w$ as a single block index gives $B^{i_0}=A^w$.
    \begin{enumerate}
        \item If $A^w$ is invertible, then
              $S_{D^2}(B)=\MN{D}$, whence
              $S_{D^2n}(A)=\MN{D}$.
        \item If $A^w$ is non-invertible, then the eigenvector $\varphi$
              spans a proper subspace, and the rank-one extraction argument
              gives $S_{D^2}(B)=\MN{D}$, whence
              $S_{D^2n}(A)=\MN{D}$.
    \end{enumerate}
    In the invertible case, the invertible Kraus operator $B^{i_0}$ allows
    direct dimension growth, giving
    $S_{D^2-k_B+1}(B)=\MN{D}$, where $k_B=\krausrk(B)$.
    In the non-invertible case, the bound $S_{D^2}(B)=\MN{D}$ follows from
    tracking indices through the rank-one extraction argument of
    \cite[Lemma~2(b)]{Sanz2010Quantum}; the separate blocked support
    construction in Theorem~\ref{thm:wielandt_lemma2b} does not supply this
    explicit bound. In both cases,
    $\iota(A)\le D^2n\le D^2(D^2-d'+1)$.
\end{proof}

\begin{remark}[Quantitative and qualitative Wielandt bounds]%
    \leanok
    The quantitative bound
    $\iota(A)\le(D^2-d'+1)D^2$
    is cited from~\cite[Theorem~1, case~(1)]{Sanz2010Quantum}.
    The later Fundamental Theorem needs only the qualitative conclusion
    $S_N(A)=\MN{D}$ for some $N$. The blocked construction in
    Theorem~\ref{thm:wielandt_lemma2b} supplies this conclusion directly from
    normality; the explicit index bound is not needed there.
\end{remark}
